\begin{figure}[htbp]
\centering
\begin{tikzcd}[column sep=5em,row sep=4em,
 every label/.append style={font=\normalsize},
 every cell/.append style={font=\large}]
 X^I \arrow[r,shift left=1ex,"\sigma"] \arrow[d,"q_I"']
 & T_{\mathcal U_X}^{*} \arrow[l,shift left=1ex,"\pi"] \arrow[d,"q_*"] \\
 \Pi_1^q(X) \arrow[r,shift left=1ex,"\overline\sigma"]
 & G_{\mathcal U_X}^{*} \arrow[l,shift left=1ex,"\overline\pi"]
\end{tikzcd}

\smallskip
\(\displaystyle
 \pi\sigma=\mathrm{id}_{X^I},\qquad
 G_{\mathcal U_X}^{\mathrm{tr}}\cong\Pi_1^q(X)
 \cong G_{\mathcal U_X}^{\mathrm{obs}}.
\)
\caption{The universal one-letter section supplies the topology, and the
universal rules supply the equivalence relation.  Here
$*\in\{\mathrm{tr},\mathrm{obs}\}$ and
$\sigma(\gamma)=(\langle\gamma\rangle,\gamma)$; $q_I$ is the homotopy
quotient map.  The lower maps are induced by $\sigma$ and $\pi$.
The upper maps are continuous, with $\pi\sigma=\mathrm{id}$, so $\pi$ is
quotient.  Agreement of scoped and homotopy relations makes the lower maps
mutually inverse homeomorphisms.  Under these identifications, $J_{\mathcal U_X}$
is the identity on path classes.}
\label{fig:universal-section}
\end{figure}
